\chapter{Paid Acquisition: How the Platforms Actually Work}\label{ch:paid-acquisition-platforms}

% Scope: Meta & Google auction/delivery, campaign structure, bidding strategies,
%   the learning phase, budget pacing, and Advantage+ / Performance Max automation.
%   Reference implementation for the digital part (cf. \cref{ch:corporate-finance-core}).

Two firms --- Google and Meta --- arbitrate most of the world's paid demand, and they do it the same way: a real-time auction that does \emph{not} simply sell to the highest bid. Understanding the value function they actually maximise tells you why a better ad is cheaper, why a tighter cost target buys fewer customers, and why fragmenting budgets quietly destroys performance.

\section{The Ad Auction}\label{sec:ad-auction}

You can control one number per auction --- your bid --- yet pay less than you bid and still lose to a lower bidder. What the platform actually ranks on, and how to move it in your favour, is the central mechanic of paid acquisition.

The platform sells attention but lives on retention: an irrelevant ad that wins on price drives users away, so both networks discount your bid by how good the ad is. Bid buys you into the auction; \emph{predicted quality} decides what that bid is worth. Raising quality is usually cheaper than raising the bid.

Meta ranks each impression by a total value that multiplies the bid by the predicted action rate and adds an experience term \parencite{metadev-advantage-campaigns}:
\begin{equation}\label{eq:metavalue}
  V_{\text{Meta}} \;=\; \underbrace{b\,\hat p}_{\text{expected revenue}} \;+\; u,
\end{equation}
where $b$ is the bid, $\hat p$ the estimated action (conversion) rate, and $u$ an ad-quality/user-value term. Google ranks by Ad Rank, the bid scaled by a Quality Score $Q$ (expected CTR, ad relevance, landing-page experience), subject to thresholds \parencite{gads1722122,gads7634668}:
\begin{equation}\label{eq:adrank}
  \mathrm{AdRank} \;=\; b \times Q \;+\;(\text{format \& context terms}).
\end{equation}
Because the auction clears at roughly the rank needed to beat the advertiser below you, your effective cost per click falls as your own quality rises:
\begin{equation}\label{eq:ecpc}
  \mathrm{CPC} \;\approx\; \frac{\mathrm{AdRank}_{\text{below}}}{Q} + \$0.01 .
\end{equation}
The mechanism is the generalised second-price auction with a quality weight --- the same structure analysed under auction theory in \cref{ch:microeconomics}. \Cref{eq:adrank,eq:metavalue} are the platforms' published skeletons; the live systems add dozens of context features and recompute $\hat p$ per user.

\begin{example}
A competitor below you clears at $\mathrm{AdRank}=16$. With a Quality Score of $8$, \cref{eq:ecpc} gives a click price of $16/8+\$0.01\approx\money{2}$. Lift relevance and landing-page experience until $Q=10$ and the same position costs $16/10+\$0.01\approx\money{1}.61$ --- a \qty{20}{\percent} cost cut bought with creative, not budget.
\end{example}

\begin{admonition}{Warning}
Treating bid as the throttle misfires. Doubling the bid on a low-quality ad can leave Ad Rank below the serving threshold, so spend rises with no extra delivery; meanwhile a rival with half your bid and twice your $Q$ outranks you. A strong brand earns higher organic CTR and thus a higher $Q$ on its own terms (\cref{ch:brand-growth}).
\end{admonition}

Both networks express this auction through a three-level hierarchy, but they put the targeting in different places (\cref{fig:campaign-hierarchy}): Google organises around \emph{keywords/intent} at the ad-group level \parencite{gads14752782,gads2375404}, Meta around \emph{audiences} at the ad-set level. Each targeting object is the operational form of the segment chosen in \cref{ch:segmentation-targeting-positioning}.

\begin{figure}[htbp]
  \centering
  \includegraphics[width=0.92\linewidth]{campaign-hierarchy}
  \caption{The tripartite campaign hierarchy. Google centres targeting on keywords (intent) at the ad-group level; Meta centres it on audiences at the ad-set level. Modern guidance consolidates both to feed the learning phase (\cref{sec:learning-phase}).}
  \label{fig:campaign-hierarchy}
\end{figure}

\section{Bidding: Telling the Machine Your Goal}\label{sec:bidding}

You no longer set per-click bids; you hand the algorithm a goal. The choice of goal --- volume, a cost target, or a return target --- and the level at which you set it determines whether delivery stabilises or profit leaks.

A bid strategy is a constraint on a maximiser. ``Maximise conversions'' spends the whole budget for the most events; adding a Target CPA or Target ROAS tells the system to walk away from auctions whose predicted cost is too high. Every constraint you add trades volume for control.

Google's Smart Bidding exposes Maximize Conversions (optionally bounded by \emph{Target CPA}) and Maximize Conversion Value (optionally bounded by \emph{Target ROAS}) \parencite{gads7381968,gads6268632,gads6268637}. Meta's Marketing API exposes the parallel set \parencite{metadev-advantage-campaigns}: \texttt{LOWEST\_COST\_WITHOUT\_CAP} (Highest Volume), \texttt{COST\_CAP}, \texttt{LOWEST\_COST\_WITH\_BID\_CAP} (a hard per-auction ceiling), and \texttt{LOWEST\_COST\_WITH\_MIN\_ROAS}. The practical regime: start on volume or lowest-cost while conversion data is thin, then add a CPA or ROAS bound once events are dense enough to optimise against.

The two efficiency metrics the goals enforce are simply
\begin{equation}\label{eq:roas-cpa}
  \mathrm{ROAS}=\frac{\text{conversion value}}{\text{ad spend}},
  \qquad
  \mathrm{CPA}=\frac{\text{ad spend}}{\text{conversions}} .
\end{equation}

\begin{example}
Our running subscription business earns \qty{42}{\percent} margin on a \money{50}/month plan: \money{21}/month, or \money{252}/year of contribution. Discounting at \qty{11}{\percent} against \qty{3}{\percent} margin growth, its customer lifetime value is the growing perpetuity of \cref{eq:gordon}:
\[
  \mathrm{LTV}=\frac{\money{252}}{0.11-0.03}=\money{3150}.
\]
Holding the textbook $\text{LTV}:\text{CAC}=3$ floor (\cref{ch:customer-economics}) caps acquisition at $\mathrm{CAC}\le\money{1050}$. So a Target CPA of \money{600} is safe (LTV:CAC $=5.25$); at \money{120000}/month that buys $120000/600=200$ conversions --- about \num{50} a week, clearing the learning-phase minimum below.
\end{example}

\begin{admonition}{Warning}
Setting the cost target at the margin instead of below it is the most common bidding error. Push Target CPA down to \money{300} ``to be efficient'' and the system abstains from most auctions: weekly events fall below \num{50}, the ad set never stabilises (\cref{sec:learning-phase}), and realised CPA \emph{rises}. The profitable move is a CPA comfortably under the LTV ceiling, not at it.
\end{admonition}

\begin{figure}[htbp]
  \centering
  \includegraphics[width=0.86\linewidth]{bid_response}
  \caption{The volume--efficiency trade-off. Tightening Target CPA lifts efficiency but cuts conversions; below the learning-phase floor (\(\sim\)50 events/week) delivery destabilises. The usable window sits between that floor and the LTV:CAC ceiling.}
  \label{fig:bid-response}
\end{figure}

\section{The Learning Phase}\label{sec:learning-phase}

The delivery model is a predictor of $\hat p$; with little data its estimates are noisy, so early cost-per-result swings wildly. It needs a minimum harvest of recent conversions before its predictions --- and your costs --- settle.

A Meta ad set exits learning after roughly \num{50} optimisation events in \num{7} days; too few and it is flagged \emph{Learning Limited} \parencite{metadev-advantage-campaigns}. Google's Smart Bidding calibrates over about \num{50} conversions or two-to-three conversion cycles \parencite{gads13020501}. A ``significant edit'' --- a budget swing past \(\sim\)\qty{20}{\percent}, a new bid strategy, or a targeting change --- resets the clock on both.

\begin{admonition}{Warning}
Fragmenting spend defeats the learning phase. Ten ad sets at \money{3000}/month each clear far fewer than 50 events apiece and \emph{all} sit in Learning Limited, splitting conversion signal and competing in the same auction. Consolidating into one or two thicker ad sets is the single highest-leverage fix --- the same throughput logic as pooling a queue (\cref{ch:operations}).
\end{admonition}

\section{Budget Pacing}\label{sec:pacing}

Neither platform divides the daily budget by 24. Google may spend up to \num{2}$\times$ the daily budget on a high-opportunity day, bounded by a monthly cap of \num{30.4}$\times$ the daily figure \parencite{gads13685469,gads10486637}. Meta now permits up to \num{1.75}$\times$ on a given day, bounded by a rolling \num{7}$\times$ weekly limit. Pacing chases the hours where $\hat p$ is highest, which is why manual dayparting usually costs more than it saves.

\section{Automation: Advantage+ and Performance Max}\label{sec:automation}

Google Performance Max collapses Search, YouTube, Display, Discover, Gmail, and Maps into one asset-fed engine that controls bidding, placement, and creative assembly toward a CPA or ROAS goal \parencite{gads10724817,gads10724896}. Meta's Advantage+ Sales does the equivalent across Meta inventory, signalled by an \texttt{advantage\_state\_info} flag once audience, budget, and placement automation are all on \parencite{metadev-advantage-plus-campaign-}.

Automation optimises for the objective it is fed --- and only that. Hand it raw ``leads'' and it will manufacture cheap, low-quality ones; hand it verified revenue or qualified-lead value and it compounds. Whether the black box transfers advertiser surplus to the platform is genuinely disputed; what is not disputed is that its output is only as good as the conversion signal fed back to it --- the subject of \cref{ch:attribution-and-measurement}.

The auction rewards quality, the bid strategy converts a business goal into a constraint, and the learning phase demands enough signal to honour it. All three depend on measuring the right conversion and returning it to the platform; that feedback loop, and the attribution behind it, is \cref{ch:attribution-and-measurement}. The official mechanics here are documented by \textcite{gads1722122} and \textcite{metadev-advantage-campaigns}.
